\subsection{Thinker: Empathetic Voice-Agent Speech-Text Foundation Model}

JoyAI-Talker adopts a standard encoder–projector–decoder architecture, integrating the JoyAI-LLM Flash language model with a speech encoder through lightweight MLP projectors. An overview of the model architecture is presented in Figure~\ref{fig:example}.

For speech understanding, the speech encoder follows an attention-based encoder–decoder architecture that has been pre-trained on large-scale automatic speech recognition (ASR) corpora. Input audio is first converted into log Mel filter-bank features and then downsampled by a factor of 8× using stacked Conv2D blocks before being processed by the Transformer encoder. This design reduces the acoustic token rate to 12.5 Hz, significantly improving computational efficiency while preserving the semantic information required for downstream language modeling.

The language backbone is built upon JoyAI-LLM Flash~\cite{cai2026joyaillmflashadvancingmidscale}, a sparse Mixture-of-Experts (MoE) large language model comprising 48.9B total parameters, with approximately 3.28B parameters activated for each input token. Following recent large-scale MoE architectures such as DeepSeek-V3~\cite{deepseekai2025deepseekv3technicalreport} and Kimi-K2~\cite{kimiteam2026kimik2openagentic}, the backbone incorporates Multi-Head Latent Attention (MLA) together with RMSNorm, Rotary Position Embeddings (RoPE), and SwiGLU feed-forward networks.

The backbone consists of 40 Transformer layers, including one dense Transformer layer followed by 39 sparse MoE layers. Each MoE layer contains a fine-grained pool of 256 experts and adopts an auxiliary-loss-free load-balancing routing strategy. For each input token, the router dynamically activates the Top-8 routed experts together with a dedicated shared expert, achieving a favorable balance between model capacity and computational efficiency.

\subsubsection{Training Recipe}

Initialized from the JoyAI-LLM Flash pretraining phase, JoyAI-Talker undergoes a unified, multi-stage speech-text joint training paradigm across all training phases. Our motivation for initiating joint multi-modal training at the mid-training phase—rather than delaying it until the SFT phase—is twofold. First, the early foundational pretraining stages establish the model's fundamental cognitive baseline and world knowledge over massive textual corpora. Second, integrating the speech modality at this earlier stage promotes stable cross-modal alignment and prevents the catastrophic forgetting of textual knowledge, thereby improving multi-modal understanding without inducing textual degradation. 

Consequently, we implement a unified speech-text mixed training scheme across all subsequent phases, sequentially optimizing the model through four distinct stages:

\begin{itemize}
    \item \textbf{Speech-Text Joint Mid-training:} This phase aims to bridge the representation space of the speech and textual modalities while preserving the foundational cognitive capabilities of the pre-trained language backbone. We introduce a highly diversified joint corpus. For the textual modality, this comprises web crawls, reasoning-intensive code repositories, high-fidelity PDF documents, and synthetic data to balance general knowledge with logical reasoning. For the speech modality, we incorporate a wide variety of tasks including automatic speech recognition (ASR), speech translation, speech question-answering (QA), audio captioning, real-world multi-turn speech dialogues, and speech text continuation (synthesized by feeding ASR transcriptions into an LLM to generate coherent textual responses paired with the original audio). This joint pretraining is split into two speech-text joint steps: 
    \begin{itemize}
        \item \textit{Step 1 (Audio Adapter Warmup):} Both the LLM backbone and the audio encoder are frozen, and we optimize only the multi-modal audio adapter to map high-dimensional speech features into the linguistic embedding space of the LLM.
        \item \textit{Step 2 (Full-Parameter Joint Training):} We unfreeze all model parameters—including the LLM backbone, the speech encoder, and the audio adapter—for comprehensive end-to-end joint optimization.
    \end{itemize}
    Throughout this phase, speech data is formatted in a chat-style Q\&A structure, while textual data is ingested as raw, continuous text to preserve cognitive capacity.
    
    \item \textbf{Speech-Text Joint Context Extension:} Following mid-training, this phase is designed to comprehensively bolster reasoning capabilities over extensive textual documents and enhance interactive proficiency during extended, multi-turn speech dialogues. All model parameters remain fully trainable. To target general long-context scaling, the textual mixture is rebalanced by upsampling long-context mathematical and coding datasets. For the speech modality, we upsample real-world multi-turn long speech dialogues and dynamically synthesize multi-turn ASR and translation tasks by concatenating single-turn samples. To emphasize long-range dependencies while preserving short-context proficiency, the joint data mixture blends long-form text (academic papers, financial reports, etc.) and long-duration audio (podcasts, lectures, etc.) with a balanced proportion of short-sequence data.
    
    \item \textbf{Speech-Text Joint Supervised Fine-Tuning (SFT):} Serving as the cornerstone of our post-training alignment pipeline, this phase is strategically designed to expand the model's multi-modal knowledge boundaries while cultivating highly colloquial, concise, and direct speech communication. The textual SFT dataset covers code, math, STEM, and general instruction-following, where we upsample knowledge-dense and interaction-centric categories to facilitate communicative voice dialogues. The speech SFT dataset covers ASR, translation, QA, speech text continuation, audio captioning, and empathetic speech dialogues. To facilitate cross-modal alignment and elevate speech response quality, we upsample synthesized speech-to-text (S2T) data constructed by converting conversational textual SFT prompts into natural speech audio via our in-house text-to-speech model.
    
    \item \textbf{Speech-Text Joint Direct Preference Optimization (DPO):} To further refine response quality, improve instruction-following, and suppress residual hallucinations, we implement a dedicated DPO phase. The model is optimized directly on pairwise preferred and dis-preferred responses self-generated by the SFT policy. For each seed prompt (spanning general QA, reasoning, and speech inputs), the SFT policy generates multiple candidate responses, from which preferred and dis-preferred pairs are selected using a multi-faceted model judging framework. Candidates are routed to three complementary, specialized judges:
    \begin{itemize}
        \item \textit{General QA Judge:} An LLM judge evaluates mathematics, code, STEM, and general knowledge responses along correctness, logical soundness, clarity, and completeness, using reference answers as quality anchors.
        \item \textit{Instruction-Following (IF) Judge:} For prompts with explicit structural constraints (exact item counts, mandatory keywords, length limits, JSON output), we prompt an LLM judge with a constraint-compliance rubric. This rubric allocates most points to hard-constraint satisfaction (with fixed penalties per violation type) and the rest to content quality and output cleanliness, thereby raising the candidate scoring margin on constraint-violating cases.
        \item \textit{Rule-Based IF Judge:} For machine-verifiable constraints, we score candidates using deterministic, rule-based validations, ensuring full reproducibility and zero inference cost.
    \end{itemize}
    For speech inputs, each prompt carries an audio field while the candidate responses remain textual, allowing us to leverage the same QA and IF judges on the transcribed text. Crucially, rather than pooling speech-conditioned and text-conditioned candidates cross-modally—which would create severe distribution mismatches and an excessively large optimization gap due to the systematic modality gap—we adopt a \textit{speech-only strategy}. This strategy constructs preference pairs strictly within the speech-conditioned pool. This local comparison provides a highly calibrated, realistic, and learnable gradient signal, enabling the model to steadily correct its speech-pathway behavior. The resulting textual and speech preference sets are then mixed for joint DPO training.
\end{itemize}

\subsubsection{Training Details}
Our audio-language model training system is built on a highly optimized extension of the Megatron-Core framework.   Alongside Data Parallelism (DP), Tensor Parallelism (TP), Sequence Parallelism (SP), and Pipeline Parallelism (PP), we employ 8-way Expert Parallelism (EP) to scale the sparse Mixture-of-Experts backbone.   MoE operations are accelerated through DeepEP for low-latency token dispatch and combination, grouped GEMM, and fused permutation and routing.   CUDA Graph capture further reduces CPU launch overhead and stabilizes step latency.

The training sequence length is configured at 8K tokens during the Mid-training Phase, and is subsequently scaled to 64K tokens during the Context Extension and SFT Phases.   For the 64K-token training stages, we adopt Context Parallelism (CP) with THD-format attention to distribute the computational load.   Under this long-context regime, we employ a best-fit sequence packing strategy utilizing block-diagonal attention masks to concatenate multiple independent samples or speech dialogue sessions into a single training sequence.   This is powered by an audio-aware packing algorithm that jointly considers token length and audio-padding waste, preventing cross-contamination between unrelated samples while approximately doubling end-to-end throughput over standard length-only packing.   We also distribute the audio encoder across CP ranks: each rank processes its local audio slice, and a differentiable all-gather reconstructs the embeddings while preserving gradient flow.   This reduces per-rank encoding cost in proportion to the CP degree and avoids an encoder bottleneck in long-context multimodal training.

Across all joint training phases, we implement a strict loss masking strategy to optimize representation learning for target generations.   For both textual corpora and multi-turn conversational datasets, the autoregressive loss is computed solely on the Assistant's responses.   The loss on all other components is completely masked out;   this includes all user-side inputs—encompassing both textual and speech queries—as well as all tool-side responses (such as tool execution results and API replies), on which no loss is computed.   This selective gradient optimization ensures that the model's representation learning capacity is focused exclusively on generating accurate, coherent, and concise assistant responses.

\subsubsection{Empathy}


Large Speech Language Models (LSLMs) have substantially advanced speech language understanding, driving applications ranging from intelligent assistants to empathetic companions \cite{DBLP:conf/naacl/WangZLSLZLAC25}. However, effective dialogue extends beyond merely interpreting \textit{what} is said, it necessitates recognizing speaker attributes (e.g., age and gender) and discerning \textit{how} the message is delivered \cite{DBLP:conf/iclr/ChengHYLF0J0Z0025, DBLP:conf/emnlp/YanLCNYMYC25}. Non-verbal acoustic cues-such as emotional prosody and paralinguistic behaviors (e.g., laughter, sigh, and cough)-are essential for capturing these nuances, facilitating natural, trustworthy, and emotionally intelligent communication. Despite their importance, existing approaches typically model these elements in isolation, failing to reflect the tightly coupled nature of understanding, reasoning, and response generation inherent in human conversation.

To address this gap, we introduce \textbf{Persona-Adaptive Empathetic Response (PAER)} to deliver highly customized, multifaceted empathic interactions that dynamically adapt to individual users. This framework is explicitly structured as a hierarchical cognitive pipeline: audio understanding $\rightarrow$ chain-of-thought (CoT) reasoning $\rightarrow$ empathetic dialogue generation. Our main contributions are twofold:

\begin{itemize}
    \item \textbf{A Progressive Cognitive Empathy Framework:} We design a progressive task formulation that mirrors human cognitive processes. Initially, we employ foundational perception tasks to accurately extract speaker attributes (e.g., gender, age, and emotional state) from the input audio. \textit{Thinker} module then explicitly integrates these attributes into its reasoning process via CoT. This facilitates the generation of context-adaptive responses that contain not
only contextually appropriate text but also precise controls over
utterance-level expression and localized paralinguistic events (e.g., sighing,
slow speaking rate, and low volume), without altering the assistant's speaker
identity. Subsequently, the Talker module, following the JoyVoice speech generation
framework~\cite{joyvoice} and refined through instruction-aware supervised
fine-tuning (SFT), translates the Thinker's textual and control outputs into
highly expressive and natural speech.
    
    \item \textbf{Comprehensive Evaluation and Superior Performance:} We rigorously evaluate our model across foundational audio understanding and downstream empathy generation tasks. First, we validate its precise perception of speaker attributes on public benchmarks, achieving outstanding performance on AIR-Bench \cite{DBLP:conf/acl/YangXLC0ZLLZZZ24} (for age and gender) and MER2025 \cite{DBLP:conf/mm/00040X0L000CZ0P25} (for emotion). These robust perceptual capabilities lay a solid foundation for individual-specific empathy. Building upon this, we comprehensively assess end-to-end empathic expressiveness using the open-source EchoMind dataset \cite{DBLP:journals/corr/abs-2510-22758}. Extensive experiments demonstrate that our approach consistently achieves superior performance in both textual and acoustic empathy metrics.
\end{itemize}

